%!TeX root=./tese.tex
% Validation-anchor chapter. Artifacts: simulation/results/live51/
% (offline twins + live_runs/{live_fc,live_re,verify_flag}); harness:
% simulation/live51/. Ledger: "5.1 validation anchor" entries.
% Upstream fix: ai-dynamo/dynamo branch fix/planner-live-fpm-observation-mode
% commit 33a9babe01 (local).
\chapter{Validation Anchor: the Loop on the Live Stack}
\label{cap:validation}
\chaptermark{Validation Anchor}

Every closed-loop result so far ran through Dynamo's \emph{offline}
replay: real planner code, virtual clock, in-process actuation. This
chapter anchors that platform to the \emph{live} stack: the same
policies, on the same workload, executed through Dynamo's real control
plane in real time --- HTTP frontend with in-process KV router,
independent mocker workers discovered over etcd, Prometheus metrics,
the SLA planner as its own service, and scaling actuated through the
planner's VirtualConnector protocol by a worker-manager written for
this thesis. The
question is the platform-validity one (RQ3), asked at the last layer
simulation cannot reach: \emph{does the simulated control loop behave
the same on the real stack?}

\section{Setup}
\label{sec:val-setup}

The anchor window is burst3, the one-minute-onset episode and the
hardest case in the fleet matrix, replayed in real time
(\SI{27}{\minute}, 3{,}093 requests) by a fixed-schedule load client
(aiperf). The client synthesizes each request at its trace timestamp with its
exact token counts and prefix-sharing structure, and measures TTFT/ITL
client-side. Both anchor policies run the production planner service:
the \emph{forecast} policy is the thesis's winning configuration
(archived Chronos p90 path at 15-second bins through the
\texttt{forecast\_file} predictor, tail-aware conversion, 60-second
replica lead); the \emph{reactive} policy is the variant-B baseline
(ARIMA predictor, 200/20\,ms window-mean targets). Engine timing is
identical to the offline runs by construction: live mocker workers load
the same AIC latency model, engine arguments, and KV-block capacity ---
one TP2 replica per OS process, so replicas die independently. The
cold-start lead is modeled by delayed spawn (59\,s + a measured
1.0\,s spawn-to-registered latency); mocker-native
\texttt{startup\_time} could not be reused because its live semantics
\emph{invert} the offline ones (Section~\ref{sec:val-findings}). Each
live run is compared against an offline twin --- identical trace,
planner configuration, capacity mapping, lead, and budget ---
re-run through the replay platform of the preceding chapters.

\section{Results: verdicts, tails, and decisions agree}
\label{sec:val-results}

\begin{table}[htbp]
\centering
\caption{Live stack versus offline twin on burst3 (evaluation SLA:
TTFT p99 $<$ \SI{1}{\second}, ITL p99 $<$
\SI{100}{\milli\second}; budget-matched at 4 TP2 replicas). ITL p99 is the
per-request inter-token latency $(\text{RequestLatency}-\text{TTFT})/(\text{OSL}-1)$
computed like-for-like in both columns, matching aiperf's live metric and
Dynamo's SLA definition (ITL${=}$TPOT).}
\label{tab:val-anchor}
\small
\begin{tabular}{@{}llrrr@{}}
\toprule
Policy & Metric & Live & Offline twin & $\Delta$ \\
\midrule
Forecast p90 & TTFT p99 (ms) & 695 & 651 & $+6.8\%$ \\
 & ITL p99 (ms) & 11.5 & 9.7 & $+1.8$\,ms \\
 & replica-seconds & 6{,}063 & 6{,}400 & $-5.3\%$ \\
 & verdict & \textbf{PASS} & \textbf{PASS} & agree \\
\midrule
Reactive & TTFT p99 (ms) & 3{,}531 & 3{,}495 & $+1.1\%$ \\
 & ITL p99 (ms) & 150.7 & 152.7 & $-1.3\%$ \\
 & replica-seconds & 1{,}799 & 1{,}948 & $-7.6\%$ \\
 & verdict & \textbf{FAIL} & \textbf{FAIL} & agree \\
\bottomrule
\end{tabular}
\end{table}

The anchor holds in both verdict directions
(Table~\ref{tab:val-anchor}): the forecast policy passes the tail SLA
live exactly as simulated; the reactive policy fails likewise, its
p99 within 1.1\% of the simulation; and provisioned
replica-time agrees within 8\%. The agreement extends below the
aggregate numbers. The \emph{demand} path is byte-identical: for the
first twelve planner ticks the live and offline planners log the same
predicted demand to the second decimal (first tick \texttt{29.86 rps} in
both, both computing \texttt{6 replicas}), because the predictor,
registry, construction order, and demand conversion are the same code;
only the clock differs. The \emph{capacity} belief, by contrast, already
differs at the first tick (live \texttt{5.71} versus offline
\texttt{5.30 engine\_rps}, both still yielding 6 replicas) and drifts
within its expected band throughout, because the regression is tuned
from sampled forward passes that arrive in different order live; but the
\emph{applied} capacity,
after the budget cap, tracks the same trajectory. The
client-observed TTFT tails land within 1.1--6.8\% of the simulated ones
(a single live run per policy, so this is a point difference between the
two systems, not a repeated-run noise band). The burst3 policy split
reproduces on the real stack in the two policies carried live: the p90
forecast passes where the reactive baseline fails (the failing p50 leg
is established offline only, Chapter~\ref{cap:closedloop}).

\section{What the live stack revealed}
\label{sec:val-findings}

The anchor's second product is the divergences it surfaced --- none
policy-level, all worth knowing before a deployment:

\begin{description}
\item[Live \texttt{startup\_time} inverts the offline semantics.]
Offline, a scaling-up mocker with \texttt{startup\_time} is invisible
to the router until ready. Live, the same option registers the worker
\emph{immediately} and stalls its requests: a routable-but-dead window
that a load-aware router, seeing an idle-looking worker,
preferentially fills; hence the delayed-spawn modeling of the Setup.
\item[The live planner's capacity model starves at calm load.] The live
planner samples worker metrics in tracking mode (the \emph{latest}
forward pass per worker, once per load tick), while the offline
replay feeds \emph{every} pass. Mixed prefill+decode passes, the only
samples that inform aggregated capacity, are rare at calm load: only
${\approx}1\%$ of calm-load samples are the mixed passes the capacity
model needs --- too few to make it ready. It would warm only once the workload was already busy,
too late for scale-ahead decisions. The fix (contributed to the Dynamo
tree as an opt-in \texttt{fpm\_observation\_mode}) drains every
published pass into the same regression feed the replay uses;
with it, the live planner is ready at its first tick and computes the
offline twin's first decision.
\item[Tick cadence stretches under stall.] The live planner's 60-second
ticks held cadence except one ${\approx}120$-second stall in the archived
run, after which the forecast feed realigned to wall time; the offline virtual clock can
never stall. Decision \emph{count} is therefore not exactly
reproducible live; decision \emph{values} at matching times are.
\item[Budget units differ between paths.] The offline replay caps
replicas at \texttt{max\_gpu\_budget}; the live planner divides the
same field by GPUs-per-replica. The anchor comparison equalizes the
caps. The cap \emph{does} bind in these episode replays --- the
tail-aware forecast requests up to 24 language replicas and
\texttt{max\_gpu\_budget}${=}8$ clamps it to 8 --- so it reshapes the
provisioned count, not the pass/fail verdict; the unit mismatch is
equalized to the same applied cap and remains a reportable upstream defect.
\end{description}

\section{Scope of the anchor}
\label{sec:val-scope}

The anchor covers one window (the hardest), two policies (the
contribution and the baseline it beats), the language pool (the pool
with a production control path; the encoder pool's controller exists
only in simulation), and mocked engine execution. It validates the
\emph{control loop}, not GPU silicon, which is what
Chapter~\ref{cap:aiconfigurator} anchored physically on the encoder
side. Within that scope the conclusion is direct: the platform's
closed-loop behavior transfers to the real stack in both verdict
directions --- the SLA verdicts are identical on all three metrics,
TTFT and provisioned replica-time agree within single-digit percent, and the
forecast policy's ITL differs by \SI{1.8}{\milli\second} (18\%, on a
\SI{9.7}{\milli\second} base, far under the \SI{100}{\milli\second} bound;
Table~\ref{tab:val-anchor}) --- and the divergences that
do exist
(identified, explained, in one case fixed on a local branch to be proposed upstream) are properties
of the harnesses, not of the policies this thesis compares.
